% TOP_PROBLEM: {"id": 2693, "title": "Kirby Problem 1.34", "category_id": 7, "category": {"id": 7, "name": "topology", "display_name": "Topology", "description": "Properties preserved under continuous deformations.", "slug": "topology", "order_index": 7, "created_at": "2026-07-31T15:26:25.671Z"}, "status": "open"}
% FIRST_POSTED: null
% ENTRY_KIND: statement_only
% Imported from: batch10.tex; UnsolvedMath ID 2693
\documentclass[11pt]{article}
\usepackage[margin=25mm]{geometry}
\usepackage{fontspec}\setmainfont{Latin Modern Roman}
\usepackage{amsmath,amssymb,mathtools}
\usepackage{unicode-math}\setmathfont{Latin Modern Math}
\usepackage{xurl,hyperref}
\hypersetup{colorlinks=true,urlcolor=blue}
\setcounter{secnumdepth}{0}
\setlength{\parindent}{0pt}\setlength{\parskip}{5pt}
\setlength{\emergencystretch}{3em}
\newcommand{\sref}[2]{\hyperlink{source:#1:#2}{[S#2]}}
\newcommand{\cataloguescope}[1]{\paragraph{Recorded target.}#1}
\begin{document}
% UnsolvedMath ID 2693; source catalogue code KP-1.34
\setcounter{section}{2692}
\section{Kirby Problem 1.34}\label{2693}
{\small\sffamily UnsolvedMath ID 2693 · Topology}
\subsection{Definitions and mathematical statement}

\paragraph{Shared notation.}$\mathbb N=\{1,2,\ldots\}$, and $\mathbb Z,\mathbb Q,\mathbb R,\mathbb C$ denote the integers, rationals, reals and complex numbers. Any other conventions are specified locally.


All knots and links below are smooth in the oriented $3$-sphere, up to ambient isotopy, unless another ambient manifold is specified. Triply graded HOMFLYPT homology refers to the Khovanov--Rozansky braid-defined theory whose graded Euler characteristic is the HOMFLYPT polynomial. Its unreduced, middle, reduced and totally reduced variants differ by basepoint reduction; totally reduced means reduction at a basepoint on every component. Functoriality means assigning maps to oriented link cobordisms, invariant under the appropriate equivalences and compatible with identity cobordisms and composition; for reduced variants the cobordisms carry paths joining the basepoints. Cautis' theories here mean the braid-defined HOMFLYPT categorifications colored by arbitrary partitions and the finite-dimensional $\mathfrak{sl}(N)$ categorifications colored by symmetric powers, identified in the original remarks.
\paragraph{Statement recorded in the supplied source.}
(a) Construct a description of Khovanov--Rozansky triply graded HOMFLYPT homology from arbitrary planar link diagrams, invariant under all Reidemeister moves, and determine whether some flavor admits functorial link-cobordism maps. (b) Obtain the same planar-diagram description and determine the same functoriality question for the link homologies of Cautis specified above. \sref{2693}{1}\sref{2693}{2}
\subsection{Short English statement}
Replace the braid-only definitions of the specified HOMFLYPT and Cautis theories by arbitrary planar diagrams and determine which variants extend to cobordism functors.
\subsection{Sources}
\begin{description}
\item[\hypertarget{source:2693:1}{S1}] ULAM.AI and the UnsolvedMath Contributors, \emph{UnsolvedMath: A Curated Collection of Open Mathematics Problems}, record 2693 (KP-1.34). CC BY 4.0. \url{https://huggingface.co/datasets/ulamai/UnsolvedMath}.
\item[\hypertarget{source:2693:2}{S2}] R. \.{I}nan\c{c} Baykur, Robion C. Kirby and Daniel Ruberman, K3 -- A New Problem List in Low-Dimensional Topology, author preliminary edition (2026), Problem 1.34 and its remarks. \url{https://math.berkeley.edu/sites/default/files/surv-295-ruberman-watermarked-author-pdf.pdf}.
\end{description}
\label{end:2693}
\end{document}
